\section{Assumption Check for Gap Bounds}\label{sec:verification}

Table~\ref{tab:assumption-bundles} already separates the structural bundles from the downstream bundles. This short section records the extra ingredients used only for quantitative transport and finite-sample certification.

The first point is conceptual: the Preservation Stack and Recompression Stack are structural. They ask whether local validity composes to the root and whether extra clean-up passes are safe. The gap theorems add a different kind of ingredient: they need a way to turn oracle distortion into objective drift, and they need empirical envelopes that connect a sampled audit to the unobserved target quantity.

The quantitative preference-gap bounds (Theorem~\ref{thm:unified-gap} and Theorem~\ref{thm:e2e}) therefore require three extra classes of assumptions:
\begin{enumerate}[leftmargin=1.5em]
    \item \textbf{Realization-to-local-law assumptions.} In the neural-operator route, external equation-(6) approximation bounds provide an upstream tolerance on compact realized calls, and Proposition~\ref{prop:neural-operator-bridge} requires explicit transfer moduli converting that tolerance into C1/C2/C3 budgets. When no such analytic route is asserted, the sampled audit estimates the realized budgets directly.
    \item \textbf{Transport assumptions.} A method-level Lipschitz/integrability condition linking expected loss drift to oracle distance after the local-law budget has been obtained. For DPO this reduces to a policy-Lipschitz condition. For GRPO-PL the formal sufficient route uses a finite Plackett--Luce fixed-ranker condition; for GRPO-RL it uses a finite pointwise Lipschitz condition for the clipped-surrogate loss. Random-utility modeling \citep{Schervish1995} is one way to motivate those assumptions, not an extra theorem by itself.
    \item \textbf{Calibration and sampling assumptions.} The accessible judge must be well enough calibrated to the target oracle on a representative held-out set, and the sampling design must satisfy the logged-propensity and positivity conditions used by Horvitz--Thompson estimation.
\end{enumerate}

The same distinction applies to the classical-literature bridge. The
Lean artifact checks the state algebra that C-TreePO actually consumes:
ordered list-homomorphism scheduling, Gray-style algebraic aggregate
state, MUD tree invariance and seedwise public-random bookkeeping, HLL
max-register merge, Count-Min additive merge, and Agarwal state-level
tree-error transport for interval, range, and width summaries. The
remaining literature facts are not hidden. HLL's ideal independent
hash source and Mellin/depoissonization estimates, Feldman's randomized
communication lower bounds, and Agarwal's compact randomized
discrepancy and geometric size constructions are exposed as typed
theorem obligations. The paper uses those facts only through their
algebraic or error-transport consequences, which are the checked Lean
interfaces listed in Appendix~\ref{app:proof-artifacts}.

These are not hidden premises. Each has a direct falsification route: approximation/transfer residual checks for the operator layer, surface-form perturbation tests for factorization/measurability, context-swap checks for context compatibility, calibration-set audits for judge mismatch, and logged-propensity diagnostics for the sampling layer.

For preferences, the application contributes a bundle of premises rather
than a new theorem. DPO supplies oracle-measurable current/reference
policies and an oracle-indexed pair generator; GRPO-PL supplies the
policy/ranker/group-generator bundle; GRPO-RL supplies current/old/reference
policy measurability, reward measurability, and an oracle-indexed group
generator. If preferences are noisy or only approximately oracle-indexed,
the theorem stack routes the residual through the same oracle-measurement
term used in the two-stage and TreePO certificates.

Readers who prefer not to rely on these extra ingredients can still use the structural theorems directly, or treat the transport and calibration constants as explicit sensitivity parameters rather than as fixed truths.
